\subsubsection{Máquina de Atwood vía formulación Lagrangiana}

Dos bloques de masas $m_1$ y $m_2$ están conectados por una cuerda ideal (inextensible, de masa despreciable) que pasa sin fricción sobre una polea fija de radio $R$ y masa despreciable, como se muestra en la Figura \ref{fig:c5_e01_p03_atwood}. Sea $x$ la distancia que ha descendido el bloque $1$ desde su posición inicial.

\begin{enumerate}
    \item Identifique el vínculo holonómico entre las posiciones de los dos bloques y justifique que el sistema tiene un solo grado de libertad.
    \item Construya el Lagrangiano en términos de la coordenada generalizada $x$.
    \item Obtenga la ecuación de movimiento y compárela con el resultado clásico obtenido por métodos newtonianos.
\end{enumerate}

\begin{figure}[!ht]
    \centering
    \begin{tikzpicture}[>=Latex, scale=1.4]
        % --- Parámetros ---
        \def\R{0.6}      % Radio de la polea (ideal, sin masa)
        \def\hUno{2.3}   % Altura inicial del bloque 1
        \def\hDos{1.3}   % Altura inicial del bloque 2

        % --- Soporte fijo del eje ---
        \coordinate (O) at (0, 2.5);
        \draw[pattern=north east lines, pattern color=gray!60] (-0.3,3.0) rectangle (0.3,3.3);
        \draw[thick] (-0.35,3.0) -- (0.35,3.0);
        \draw[thick] (0,3.0) -- (O);

        % --- Polea (ideal) ---
        \draw[thick] (O) circle (\R);

        % --- Cuerda sobre la polea ---
        \draw[thick] ({-\R}, {2.5}) arc (180:0:\R);

        % --- Cuerda y bloque 1 ---
        \draw[thick] (-\R, 2.5) -- (-\R, {2.5-\hUno+0.4});
        \draw[thick, fill=gray!20] (-\R-0.4, {2.5-\hUno}) rectangle (-\R+0.4, {2.5-\hUno+0.4});
        \node at (-\R, {2.5-\hUno+0.2}) {$m_1$};

        % --- Cuerda y bloque 2 ---
        \draw[thick] (\R, 2.5) -- (\R, {2.5-\hDos+0.4});
        \draw[thick, fill=gray!20] (\R-0.4, {2.5-\hDos}) rectangle (\R+0.4, {2.5-\hDos+0.4});
        \node at (\R, {2.5-\hDos+0.2}) {$m_2$};

        % --- Coordenada x (desplazamiento de m1) ---
        \draw[->, thick, azul] (-\R-0.8, {2.5-\hUno+0.4}) -- (-\R-0.8, {2.5-\hUno-0.7});
        \node[left, color=azul] at (-\R-0.8, {2.5-\hUno-0.15}) {$x$};

    \end{tikzpicture}
    \caption{Máquina de Atwood: dos masas conectadas por una cuerda inextensible sobre una polea ideal.}
    \label{fig:c5_e01_p03_atwood}
\end{figure}
